\documentclass{article}
\usepackage[spanish]{babel}
\usepackage[a4paper,
            top=2cm,
            bottom=2cm,
            left=3cm,
            right=3cm,
            headheight=36pt,
            nomarginpar,
            includehead,
            includefoot,
            ]{geometry}
\usepackage{graphicx}
\usepackage{listings}
\usepackage{fancyhdr}
\usepackage{parskip}
\usepackage{courier}
\usepackage{xcolor}
\usepackage{minted}


% \usemintedstyle{monokai}
\definecolor{background}{gray}{0.2}


\setminted[python3]{
    style=default,
    % bgcolor=background,
    breaklines,
    linenos=true,
    firstnumber=0,
    fontsize=\footnotesize,
    frame=single,
    samepage=true,
    % xleftmargin=20pt,
    % xrightmargin=20pt,
}


\lstset{
    language=Python,                    % the language of the code
    basicstyle=\footnotesize\ttfamily,  % the size of the fonts that are used for the code
    breakatwhitespace=true,             % sets if automatic breaks should only happen at whitespace
    breaklines=true,                    % sets automatic line breaking
    keepspaces=true,                    % keeps spaces in text, useful for keeping indentation of code (possibly needs columns=flexible)
    escapeinside={\%*}{*)},             % if you want to add LaTeX within your code
    extendedchars=true,                 % lets you use non-ASCII characters; for 8-bits encodings only, does not work with UTF-8
    otherkeywords={None,Entrada,Salida},% if you want to add more keywords to the set
    deletekeywords={del},               % if you want to delete keywords from the given language
    literate=   {á}{{\'a}}1             % fixes utf-8 errors in lstlisting environments
                {é}{{\'e}}1
                {í}{{\'i}}1
                {ó}{{\'o}}1
                {ú}{{\'u}}1,
}


% header & footer settings
\fancyhf{}
\lhead{\includegraphics[height=32pt]{img/logo-uncuyo-fing.pdf}}
\rhead{ Licenciatura en Ciencias de la Computación \\
        Algoritmos y Estructuras de Datos II \\
        Árboles Balanceados: AVL}
\rfoot{\thepage}
\pagestyle{fancy}


\title{Trabajo Práctico Nº 1}
\author{Gonzalo Padilla}
\date{Marzo 2024}


\begin{document}
\maketitle
\pagestyle{fancy}


\section*{Parte 1}

\textbf{Importante:} Los ejercicios de esta primera parte tienen como objetivo codificar las diferentes funciones básicas necesarias para implementar un árbol AVL.

A partir de estructuras definidas como:

\begin{lstlisting}
class AVLTree:
    root = None

class AVLNode:
    parent = None
    leftnode = None
    rightnode = None
    key = None
    value = None
    bf = None
\end{lstlisting}

Copiar y adaptar todas las operaciones del \textbf{binarytree.py} (i.e. insert(), delete(), search(), etc.) al nuevo módulo \textbf{avltree.py}. Notar que estos luego deberán ser implementados para cumplir que la propiedad de un árbol AVL
\subsubsection*{Solución}
Ver ejercicios 4 y 5.


\subsection*{Ejercicio 1}
Crear un modulo de nombre \textbf{avltree.py}. Implementar las siguientes funciones:
\begin{lstlisting}
%*\textbf{rotateLeft(Tree,avlnode)}*)
    %*\textbf{Descripción}*): Implementa la operación rotación a la izquierda
    Entrada: Un Tree junto a un AVLnode sobre el cual se va a operar la rotación a la izquierda
    Salida: retorna la nueva raíz

%*\textbf{rotateRight(Tree,avlnode)}*)
    %*\textbf{Descripción}*): Implementa la operación rotación a la derecha
    Entrada: Un Tree junto a un AVLnode sobre el cual se va a operar la rotación a la derecha
    Salida: retorna la nueva raíz
\end{lstlisting}
\subsubsection*{Solución}
\inputminted{python3}{./code/snippets/ejercicio1.py}


\subsection*{Ejercicio 2}
Implementar una función recursiva que calcule el elemento balanceFactor de cada subárbol siguiendo la siguiente especificación:
\begin{lstlisting}
%*\textbf{calculateBalance(AVLTree)}*)
    %*\textbf{Descripción}*): Calcula el factor de balanceo de un árbol binario de búsqueda.
    Entrada: El árbol AVL sobre el cual se quiere operar.
    Salida: El árbol AVL con el valor de balanceFactor para cada subarbol
\end{lstlisting}
\subsubsection*{Solución}
\inputminted{python3}{./code/snippets/ejercicio2.py}


\subsection*{Ejercicio 3}
Implementar una funcion en el modulo avltree.py de acuerdo a  las siguientes especifcaciones:

\begin{lstlisting}
%*\textbf{reBalance(AVLTree)}*)
    %*\textbf{Descripción}*): balancea un árbol binario de búsqueda. Para esto se deberá primero calcular el balanceFactor del árbol y luego en función de esto aplicar la estrategia de rotación que corresponda.
    Entrada: El árbol binario de tipo AVL  sobre el cual se quiere operar.
    Salida: Un árbol binario de búsqueda balanceado. Es decir luego de esta operación se cumple que la altura (h) de su subárbol derecho e izquierdo difieren a lo sumo en una unidad.
\end{lstlisting}
\subsubsection*{Solución}
\setlength{\fboxsep}{0pt}
\inputminted{python3}{./code/snippets/ejercicio3.py}


\subsection*{Ejercicio 4}
Implementar la operación \textbf{insert()} en  el módulo \textbf{avltree.py} garantizando que el árbol  binario resultante sea un árbol AVL. 
\subsubsection*{Solución}
\inputminted{python3}{./code/snippets/ejercicio4.py}


\subsection*{Ejercicio 5}
Implementar la operación \textbf{delete()} en  el módulo \textbf{avltree.py} garantizando que el árbol  binario resultante sea un árbol AVL.
\subsubsection*{Solución}
\inputminted{python3}{./code/snippets/ejercicio5.py}


\section*{Parte 2}
\subsection*{Ejercicio 6}
\begin{enumerate}
\item Responder V o F y justificar su respuesta:
\begin{enumerate}
\item \underline{ F } En un AVL el penúltimo nivel tiene que estar completo
\item \underline{ V } Un AVL donde todos los nodos tengan factor de balance 0 es completo
\item \underline{ V } En la inserción en un AVL, si al actualizarle el factor de balance al padre del nodo insertado éste no se desbalanceó, entonces no hay que seguir verificando hacia arriba porque no hay cambios en los factores de balance.
\item \underline{ V } En todo AVL existe al menos un nodo con factor de balance 0.
\end{enumerate}
\end{enumerate}
\subsubsection*{Solución}


\subsection*{Ejercicio 7}
Sean $A$ y $B$ dos AVL de $m$ y $n$ nodos respectivamente y sea $x$ un key cualquiera de forma tal que para todo key $a\in A$ y para todo key $b\in B$ se cumple que $a < x < b$. Plantear un algoritmo $O(\log n + \log m)$ que devuelva un AVL que contenga los key de $A$, el key $x$ y los key de $B$.
\subsubsection*{Solución}


\subsection*{Ejercicio 8}
Considere una rama truncada en un AVL como un camino simple desde la raíz hacia un nodo que tenga una referencia None (que le falte algún hijo). Demuestre que la mínima longitud (cantidad de aristas) que puede tener una rama truncada  en un AVL de altura $h$ es $h/2$ (tomando la parte entera por abajo). 

Cualquier camino desde la raíz hasta un nodo que no esté completo puede ser una rama truncada según la definición del ejercicio. Dicho nodo puede no ser necesariamente un nodo hoja.
\subsubsection*{Solución}


\section*{Parte 3}
\subsection*{Ejercicios Opcionales}
\begin{enumerate}
    \item Si n es la cantidad de nodos en un árbol AVL, implemente la operación \textbf{height()} en el módulo \textbf{avltree.py} que determine su altura en $O(\log n)$. Justifique el por qué de dicho orden.
    \subsubsection*{Solución}
    \inputminted{python3}{./code/snippets/ejercicioopcional1.py}
    
    \item Considere una modificación en el módulo \textbf{avltree.py} donde a cada nodo se le ha agregado el campo \textbf{count} que almacena el número de nodos que hay en el subárbol en el que él es raíz. Programe un algoritmo $O(\log n)$ que determine la cantidad de nodos en el árbol cuyo valor del key se encuentra en un intervalo $[a, b]$ dado como parámetro. Explique brevemente por qué el algoritmo programado por usted tiene dicho orden.
    \subsubsection*{Solución}
\end{enumerate}


\end{document}
